
\documentclass{article}
\usepackage{colortbl}
\usepackage{makecell}
\usepackage{multirow}
\usepackage{supertabular}

\begin{document}

\newcounter{utterance}

\centering \large Interaction Transcript for game `cladder', experiment `full\_v1.5\_default', episode 2151 with score{-}Qwen3.5{-}9B{-}sp.
\vspace{24pt}

{ \footnotesize  \setcounter{utterance}{1}
\setlength{\tabcolsep}{0pt}
\begin{supertabular}{c@{$\;$}|p{.15\linewidth}@{}p{.15\linewidth}p{.15\linewidth}p{.15\linewidth}p{.15\linewidth}p{.15\linewidth}}
    \# & \multicolumn{2}{c}{Player} && \multicolumn{2}{c}{Game Master} \\
    \hline

    \theutterance \stepcounter{utterance}  
    & & & \multicolumn{4}{p{0.6\linewidth}}{
        \cellcolor[rgb]{0.9,0.9,0.9}{
            \makecell[{{p{\linewidth}}}]{
                \texttt{\tiny{[P1$\langle$GM]}}
                \texttt{You are an expert in causal inference. The following question is not a typical commonsense query, but rather a meticulously designed question created by a professor specializing in causal inference, intended to assess the students' mastery of the course content.} \\
\texttt{Imagine a self{-}contained, hypothetical world with only the following conditions, and without any unmentioned factors or causal relationships: Gender has a direct effect on department competitiveness and admission status. Residency status has a direct effect on department competitiveness and admission status. Department competitiveness has a direct effect on admission status. Residency status is unobserved. The overall probability of male gender is 13\%. For individuals who are not male, the probability of admission acceptance is 61\%. For individuals who are male, the probability of admission acceptance is 65\%. Is admission acceptance less likely than admission rejection overall?} \\
\texttt{Start your answer with 'Yes' or 'No', followed by additional reasoning or evidence to support your explanation.} \\
            }
        }
    }
     \\ \\

    \theutterance \stepcounter{utterance}  
    & \multicolumn{4}{p{0.6\linewidth}}{
        \cellcolor[rgb]{0.561,0.737,0.561}{
            \makecell[{{p{\linewidth}}}]{
                \texttt{\tiny{[P1$\rangle$GM]}}
                \texttt{No} \\
\\ 
\texttt{**Reasoning:**} \\
\\ 
\texttt{To determine whether admission acceptance is less likely than rejection overall, we must calculate the total probability of admission acceptance, \$P(\textbackslash{}text\{Admission\})\$, by summing the probabilities of admission for each gender group weighted by their respective population proportions.} \\
\\ 
\texttt{1.  **Identify the given probabilities:**} \\
\texttt{    *   Probability of being Male: \$P(\textbackslash{}text\{Male\}) = 0.13\$.} \\
\texttt{    *   Probability of being Not Male (Female): \$P(\textbackslash{}text\{Not Male\}) = 1 {-} 0.13 = 0.87\$.} \\
\texttt{    *   Probability of Admission given Not Male: \$P(\textbackslash{}text\{Admission\} | \textbackslash{}text\{Not Male\}) = 0.61\$.} \\
\texttt{    *   Probability of Admission given Male: \$P(\textbackslash{}text\{Admission\} | \textbackslash{}text\{Male\}) = 0.65\$.} \\
\\ 
\texttt{2.  **Calculate the total probability of admission using the Law of Total Probability:**} \\
\texttt{    \$\$P(\textbackslash{}text\{Admission\}) = P(\textbackslash{}text\{Admission\} | \textbackslash{}text\{Male\}) \textbackslash{}times P(\textbackslash{}text\{Male\}) + P(\textbackslash{}text\{Admission\} | \textbackslash{}text\{Not Male\}) \textbackslash{}times P(\textbackslash{}text\{Not Male\})\$\$} \\
\texttt{    \$\$P(\textbackslash{}text\{Admission\}) = (0.65 \textbackslash{}times 0.13) + (0.61 \textbackslash{}times 0.} \\
            }
        }
    }
    & & \\ \\

    \theutterance \stepcounter{utterance}  
    & & & \multicolumn{2}{p{0.3\linewidth}}{
        \cellcolor[rgb]{0.9,0.9,0.9}{
            \makecell[{{p{\linewidth}}}]{
                \texttt{\tiny{[GM$|$GM]}}
                \texttt{No} \\
            }
        }
    }
    & & \\ \\

    \theutterance \stepcounter{utterance}  
    & & & \multicolumn{2}{p{0.3\linewidth}}{
        \cellcolor[rgb]{0.9,0.9,0.9}{
            \makecell[{{p{\linewidth}}}]{
                \texttt{\tiny{[GM$|$GM]}}
                \texttt{no} \\
            }
        }
    }
    & & \\ \\

    \theutterance \stepcounter{utterance}  
    & & & \multicolumn{2}{p{0.3\linewidth}}{
        \cellcolor[rgb]{0.9,0.9,0.9}{
            \makecell[{{p{\linewidth}}}]{
                \texttt{\tiny{[GM$|$GM]}}
                \texttt{game\_result = WIN} \\
            }
        }
    }
    & & \\ \\

\end{supertabular}
}

\end{document}
